% ==============================================================================
% Plantilla Institucional de Trabajo Terminal — UPIIZ IPN
% Archivo: capitulos/protocolo/03-antecedentes.tex
% ==============================================================================
% ESTRUCTURA NORMATIVA DE PROTOCOLO: CAPÍTULO 3 - ANTECEDENTES
% ==============================================================================
% LÍMITE INSTITUCIONAL: Máximo 1 página.
%
% Requisitos normativos (Anexo 1):
% - Describir trabajos realizados anteriormente por:
%   * alumnos;
%   * asesores;
%   * o dentro de la Unidad Académica (UPIIZ-IPN),
%   relacionados directamente con el trabajo a desarrollar.
%
% Sugerencia editorial (no obligatoria):
% - Se puede hacer una breve mención al contexto o evolución del área si aporta
%   claridad técnica directa a la propuesta.
% ==============================================================================

\chapter{Antecedentes}
\label{cap:proto-antecedentes}

% [INSTRUCCIÓN PARA EL ALUMNO]:
% Desarrolle este capítulo en un máximo de una página abordando los antecedentes
% solicitados expresamente en el Anexo 1:

\section{Trabajos previos relacionados con la propuesta}

[Describa los proyectos de investigación, Trabajos Terminales, desarrollos tecnológicos, tesis o prototipos elaborados anteriormente por alumnos, por los asesores del proyecto o dentro de los laboratorios y academias de la UPIIZ-IPN que estén vinculados directamente con el trabajo a desarrollar].

\vspace{0.8em}
[Identifique las características, aportaciones o limitaciones de dichos trabajos previos, explicando de qué manera sirven como base, punto de partida o motivación para justificar la presente propuesta de desarrollo].

% Sugerencia editorial complementaria (opcional):
% Si resulta indispensable para contextualizar la propuesta, puede incluirse una breve
% referencia a la evolución tecnológica del área temática.
